\chapter{FCS：现实如何影响主体}
\chaptervoice{这一章，我们围绕“FCS：现实如何影响主体”做一件克制的事：把直觉压缩为状态、关系、时间尺度和可检验后果。它在全书中的任务，是说明认知内核通过身体进入现实。}

\section*{理论来路与依赖接口}
本章的直接思想来源是原文关于 Body Runtime、FBI、SGC/FCS、三信号通道、BPCC 与 Reality Authority 的具身讨论。我们采用原文较晚形成的定义来整理较早讨论，不把对话中的试探性比喻与最终模型并列成互相竞争的“定律”。已有理论锚点主要是具身认知、分层控制、MPC、状态估计与安全包络（模型预测控制、Kalman 状态估计与具身认知）。

本章使用上一章“SGC 与现实：从 Scene Graph 到 Actionable Reality”已经得到的对象与边界；本章产出将供下一章“SGC—FCS 二元现实模型”使用。若后文术语偶尔出现，只作为路线提示，不参与本章推导。

\section*{从哪里出发}
我们已经拥有上一章给出的概念，现在只向前走一步。我会先把对象放进闭环，再讨论它的数学形式，最后检查工程边界。读者不必预先相信本章术语；只需暂时接受一个方法论约定：智能结构的意义来自它在现实反馈中的作用，而不是来自名称本身。

令系统在观察窗口 $[t,t+T]$ 内的状态轨迹为 $x_{[t,t+T]}$，环境扰动为 $d$，行动为 $u$。本章讨论的对象若确实有效，就应改变条件分布 $p(x_{t+T}\mid x_t,u,d)$，或改变达到同一结果所需的代价。这个起点把讨论锚定在具身认知、分层控制、MPC、状态估计与安全包络上。
\section{Functional Feeling Coordinate/Field System}
\begin{narration}
现在把镜头移到“Functional Feeling Coordinate/Field System”。我们只沿本章已经取得的概念向前一步：先确定它在闭环里的角色，再检查它的尺度与失败边界。
\end{narration}

本节把“Functional Feeling Coordinate/Field System”落实为把原始身体量转换为强度、导数和置信度组成的主体影响场。这个定义不是说“任何相似现象都算”，而是要求我们同时给出对象域、允许的干预以及观察窗口。对于主题“FCS：现实如何影响主体”而言，它承担的是认知内核通过身体进入现实中的一个可辨识环节。

把这一关系写成最小数学骨架，可得

\begin{equation}f_k=(\sigma_k,\dot\sigma_k,q_k).\end{equation}

式中的符号只承担结构角色，具体参数仍需由任务和身体数据辨识。我们首先检查可观测量，再检查控制量，最后检查比它更慢的约束。这样的顺序来自具身认知、分层控制、MPC、状态估计与安全包络，但本书对Functional Feeling Coordinate/Field System的命名和组合仍是需要实验验证的建模提案。

这一小节的反例同样重要：原始传感值与功能影响混同会让不同身体无法使用统一语义。因此评价它不能只看一次任务结果，而要比较干预前后的轨迹、恢复时间、维护成本和风险。如果改变该机制而所有允许观察下的分布都不变，那么它至少在当前实验族中是冗余描述。

\begin{thought}
对“Functional Feeling Coordinate/Field System”做一次消融：冻结它、删除它，或只允许它以相邻时间尺度交换残差。哪些现象立即消失，哪些只在长期出现？若答案无法区分三种干预，我们还没有找到正确的状态变量。
\end{thought}

最后，把结论交还现实：本节只建立功能层关系，不由此推出特定脑区实现、主观意识或宇宙目的。可测状态、干预后果和明确失败条件，才是从原文直觉走向可检验理论的桥。

\section{Feeling 的工程定义}
\begin{narration}
现在把镜头移到“Feeling 的工程定义”。我们只沿本章已经取得的概念向前一步：先确定它在闭环里的角色，再检查它的尺度与失败边界。
\end{narration}

本节把“Feeling 的工程定义”落实为把原始身体量转换为强度、导数和置信度组成的主体影响场。这个定义不是说“任何相似现象都算”，而是要求我们同时给出对象域、允许的干预以及观察窗口。对于主题“FCS：现实如何影响主体”而言，它承担的是认知内核通过身体进入现实中的一个可辨识环节。

把这一关系写成最小数学骨架，可得

\begin{equation}f_k=(\sigma_k,\dot\sigma_k,q_k).\end{equation}

式中的符号只承担结构角色，具体参数仍需由任务和身体数据辨识。我们首先检查可观测量，再检查控制量，最后检查比它更慢的约束。这样的顺序来自具身认知、分层控制、MPC、状态估计与安全包络，但本书对Feeling 的工程定义的命名和组合仍是需要实验验证的建模提案。

这一小节的反例同样重要：原始传感值与功能影响混同会让不同身体无法使用统一语义。因此评价它不能只看一次任务结果，而要比较干预前后的轨迹、恢复时间、维护成本和风险。如果改变该机制而所有允许观察下的分布都不变，那么它至少在当前实验族中是冗余描述。

\begin{thought}
对“Feeling 的工程定义”做一次消融：冻结它、删除它，或只允许它以相邻时间尺度交换残差。哪些现象立即消失，哪些只在长期出现？若答案无法区分三种干预，我们还没有找到正确的状态变量。
\end{thought}

最后，把结论交还现实：本节只建立功能层关系，不由此推出特定脑区实现、主观意识或宇宙目的。可测状态、干预后果和明确失败条件，才是从原文直觉走向可检验理论的桥。

\section{Raw Value $\neq$ Feeling Value}
\begin{narration}
现在把镜头移到“Raw Value $\neq$ Feeling Value”。我们只沿本章已经取得的概念向前一步：先确定它在闭环里的角色，再检查它的尺度与失败边界。
\end{narration}

本节把“Raw Value $\neq$ Feeling Value”落实为把原始身体量转换为强度、导数和置信度组成的主体影响场。这个定义不是说“任何相似现象都算”，而是要求我们同时给出对象域、允许的干预以及观察窗口。对于主题“FCS：现实如何影响主体”而言，它承担的是认知内核通过身体进入现实中的一个可辨识环节。

把这一关系写成最小数学骨架，可得

\begin{equation}f_k=(\sigma_k,\dot\sigma_k,q_k).\end{equation}

式中的符号只承担结构角色，具体参数仍需由任务和身体数据辨识。我们首先检查可观测量，再检查控制量，最后检查比它更慢的约束。这样的顺序来自具身认知、分层控制、MPC、状态估计与安全包络，但本书对Raw Value $\neq$ Feeling Value的命名和组合仍是需要实验验证的建模提案。

这一小节的反例同样重要：原始传感值与功能影响混同会让不同身体无法使用统一语义。因此评价它不能只看一次任务结果，而要比较干预前后的轨迹、恢复时间、维护成本和风险。如果改变该机制而所有允许观察下的分布都不变，那么它至少在当前实验族中是冗余描述。

\begin{thought}
对“Raw Value $\neq$ Feeling Value”做一次消融：冻结它、删除它，或只允许它以相邻时间尺度交换残差。哪些现象立即消失，哪些只在长期出现？若答案无法区分三种干预，我们还没有找到正确的状态变量。
\end{thought}

最后，把结论交还现实：本节只建立功能层关系，不由此推出特定脑区实现、主观意识或宇宙目的。可测状态、干预后果和明确失败条件，才是从原文直觉走向可检验理论的桥。

\section{Energy Pressure}
\begin{narration}
现在把镜头移到“Energy Pressure”。我们只沿本章已经取得的概念向前一步：先确定它在闭环里的角色，再检查它的尺度与失败边界。
\end{narration}

本节把“Energy Pressure”落实为把原始身体量转换为强度、导数和置信度组成的主体影响场。这个定义不是说“任何相似现象都算”，而是要求我们同时给出对象域、允许的干预以及观察窗口。对于主题“FCS：现实如何影响主体”而言，它承担的是认知内核通过身体进入现实中的一个可辨识环节。

把这一关系写成最小数学骨架，可得

\begin{equation}f_k=(\sigma_k,\dot\sigma_k,q_k).\end{equation}

式中的符号只承担结构角色，具体参数仍需由任务和身体数据辨识。我们首先检查可观测量，再检查控制量，最后检查比它更慢的约束。这样的顺序来自具身认知、分层控制、MPC、状态估计与安全包络，但本书对Energy Pressure的命名和组合仍是需要实验验证的建模提案。

这一小节的反例同样重要：原始传感值与功能影响混同会让不同身体无法使用统一语义。因此评价它不能只看一次任务结果，而要比较干预前后的轨迹、恢复时间、维护成本和风险。如果改变该机制而所有允许观察下的分布都不变，那么它至少在当前实验族中是冗余描述。

\begin{thought}
对“Energy Pressure”做一次消融：冻结它、删除它，或只允许它以相邻时间尺度交换残差。哪些现象立即消失，哪些只在长期出现？若答案无法区分三种干预，我们还没有找到正确的状态变量。
\end{thought}

最后，把结论交还现实：本节只建立功能层关系，不由此推出特定脑区实现、主观意识或宇宙目的。可测状态、干预后果和明确失败条件，才是从原文直觉走向可检验理论的桥。

\section{Thermal Pressure}
\begin{narration}
现在把镜头移到“Thermal Pressure”。我们只沿本章已经取得的概念向前一步：先确定它在闭环里的角色，再检查它的尺度与失败边界。
\end{narration}

本节把“Thermal Pressure”落实为把原始身体量转换为强度、导数和置信度组成的主体影响场。这个定义不是说“任何相似现象都算”，而是要求我们同时给出对象域、允许的干预以及观察窗口。对于主题“FCS：现实如何影响主体”而言，它承担的是认知内核通过身体进入现实中的一个可辨识环节。

把这一关系写成最小数学骨架，可得

\begin{equation}f_k=(\sigma_k,\dot\sigma_k,q_k).\end{equation}

式中的符号只承担结构角色，具体参数仍需由任务和身体数据辨识。我们首先检查可观测量，再检查控制量，最后检查比它更慢的约束。这样的顺序来自具身认知、分层控制、MPC、状态估计与安全包络，但本书对Thermal Pressure的命名和组合仍是需要实验验证的建模提案。

这一小节的反例同样重要：原始传感值与功能影响混同会让不同身体无法使用统一语义。因此评价它不能只看一次任务结果，而要比较干预前后的轨迹、恢复时间、维护成本和风险。如果改变该机制而所有允许观察下的分布都不变，那么它至少在当前实验族中是冗余描述。

\begin{thought}
对“Thermal Pressure”做一次消融：冻结它、删除它，或只允许它以相邻时间尺度交换残差。哪些现象立即消失，哪些只在长期出现？若答案无法区分三种干预，我们还没有找到正确的状态变量。
\end{thought}

最后，把结论交还现实：本节只建立功能层关系，不由此推出特定脑区实现、主观意识或宇宙目的。可测状态、干预后果和明确失败条件，才是从原文直觉走向可检验理论的桥。

\section{Integrity Pressure}
\begin{narration}
现在把镜头移到“Integrity Pressure”。我们只沿本章已经取得的概念向前一步：先确定它在闭环里的角色，再检查它的尺度与失败边界。
\end{narration}

本节把“Integrity Pressure”落实为把原始身体量转换为强度、导数和置信度组成的主体影响场。这个定义不是说“任何相似现象都算”，而是要求我们同时给出对象域、允许的干预以及观察窗口。对于主题“FCS：现实如何影响主体”而言，它承担的是认知内核通过身体进入现实中的一个可辨识环节。

把这一关系写成最小数学骨架，可得

\begin{equation}f_k=(\sigma_k,\dot\sigma_k,q_k).\end{equation}

式中的符号只承担结构角色，具体参数仍需由任务和身体数据辨识。我们首先检查可观测量，再检查控制量，最后检查比它更慢的约束。这样的顺序来自具身认知、分层控制、MPC、状态估计与安全包络，但本书对Integrity Pressure的命名和组合仍是需要实验验证的建模提案。

这一小节的反例同样重要：原始传感值与功能影响混同会让不同身体无法使用统一语义。因此评价它不能只看一次任务结果，而要比较干预前后的轨迹、恢复时间、维护成本和风险。如果改变该机制而所有允许观察下的分布都不变，那么它至少在当前实验族中是冗余描述。

\begin{thought}
对“Integrity Pressure”做一次消融：冻结它、删除它，或只允许它以相邻时间尺度交换残差。哪些现象立即消失，哪些只在长期出现？若答案无法区分三种干预，我们还没有找到正确的状态变量。
\end{thought}

最后，把结论交还现实：本节只建立功能层关系，不由此推出特定脑区实现、主观意识或宇宙目的。可测状态、干预后果和明确失败条件，才是从原文直觉走向可检验理论的桥。

\section{Threat}
\begin{narration}
现在把镜头移到“Threat”。我们只沿本章已经取得的概念向前一步：先确定它在闭环里的角色，再检查它的尺度与失败边界。
\end{narration}

本节把“Threat”落实为把原始身体量转换为强度、导数和置信度组成的主体影响场。这个定义不是说“任何相似现象都算”，而是要求我们同时给出对象域、允许的干预以及观察窗口。对于主题“FCS：现实如何影响主体”而言，它承担的是认知内核通过身体进入现实中的一个可辨识环节。

把这一关系写成最小数学骨架，可得

\begin{equation}f_k=(\sigma_k,\dot\sigma_k,q_k).\end{equation}

式中的符号只承担结构角色，具体参数仍需由任务和身体数据辨识。我们首先检查可观测量，再检查控制量，最后检查比它更慢的约束。这样的顺序来自具身认知、分层控制、MPC、状态估计与安全包络，但本书对Threat的命名和组合仍是需要实验验证的建模提案。

这一小节的反例同样重要：原始传感值与功能影响混同会让不同身体无法使用统一语义。因此评价它不能只看一次任务结果，而要比较干预前后的轨迹、恢复时间、维护成本和风险。如果改变该机制而所有允许观察下的分布都不变，那么它至少在当前实验族中是冗余描述。

\begin{thought}
对“Threat”做一次消融：冻结它、删除它，或只允许它以相邻时间尺度交换残差。哪些现象立即消失，哪些只在长期出现？若答案无法区分三种干预，我们还没有找到正确的状态变量。
\end{thought}

最后，把结论交还现实：本节只建立功能层关系，不由此推出特定脑区实现、主观意识或宇宙目的。可测状态、干预后果和明确失败条件，才是从原文直觉走向可检验理论的桥。

\section{Arousal}
\begin{narration}
现在把镜头移到“Arousal”。我们只沿本章已经取得的概念向前一步：先确定它在闭环里的角色，再检查它的尺度与失败边界。
\end{narration}

本节把“Arousal”落实为把原始身体量转换为强度、导数和置信度组成的主体影响场。这个定义不是说“任何相似现象都算”，而是要求我们同时给出对象域、允许的干预以及观察窗口。对于主题“FCS：现实如何影响主体”而言，它承担的是认知内核通过身体进入现实中的一个可辨识环节。

把这一关系写成最小数学骨架，可得

\begin{equation}f_k=(\sigma_k,\dot\sigma_k,q_k).\end{equation}

式中的符号只承担结构角色，具体参数仍需由任务和身体数据辨识。我们首先检查可观测量，再检查控制量，最后检查比它更慢的约束。这样的顺序来自具身认知、分层控制、MPC、状态估计与安全包络，但本书对Arousal的命名和组合仍是需要实验验证的建模提案。

这一小节的反例同样重要：原始传感值与功能影响混同会让不同身体无法使用统一语义。因此评价它不能只看一次任务结果，而要比较干预前后的轨迹、恢复时间、维护成本和风险。如果改变该机制而所有允许观察下的分布都不变，那么它至少在当前实验族中是冗余描述。

\begin{thought}
对“Arousal”做一次消融：冻结它、删除它，或只允许它以相邻时间尺度交换残差。哪些现象立即消失，哪些只在长期出现？若答案无法区分三种干预，我们还没有找到正确的状态变量。
\end{thought}

最后，把结论交还现实：本节只建立功能层关系，不由此推出特定脑区实现、主观意识或宇宙目的。可测状态、干预后果和明确失败条件，才是从原文直觉走向可检验理论的桥。

\section{Sensory Load}
\begin{narration}
现在把镜头移到“Sensory Load”。我们只沿本章已经取得的概念向前一步：先确定它在闭环里的角色，再检查它的尺度与失败边界。
\end{narration}

本节把“Sensory Load”落实为把原始身体量转换为强度、导数和置信度组成的主体影响场。这个定义不是说“任何相似现象都算”，而是要求我们同时给出对象域、允许的干预以及观察窗口。对于主题“FCS：现实如何影响主体”而言，它承担的是认知内核通过身体进入现实中的一个可辨识环节。

把这一关系写成最小数学骨架，可得

\begin{equation}f_k=(\sigma_k,\dot\sigma_k,q_k).\end{equation}

式中的符号只承担结构角色，具体参数仍需由任务和身体数据辨识。我们首先检查可观测量，再检查控制量，最后检查比它更慢的约束。这样的顺序来自具身认知、分层控制、MPC、状态估计与安全包络，但本书对Sensory Load的命名和组合仍是需要实验验证的建模提案。

这一小节的反例同样重要：原始传感值与功能影响混同会让不同身体无法使用统一语义。因此评价它不能只看一次任务结果，而要比较干预前后的轨迹、恢复时间、维护成本和风险。如果改变该机制而所有允许观察下的分布都不变，那么它至少在当前实验族中是冗余描述。

\begin{thought}
对“Sensory Load”做一次消融：冻结它、删除它，或只允许它以相邻时间尺度交换残差。哪些现象立即消失，哪些只在长期出现？若答案无法区分三种干预，我们还没有找到正确的状态变量。
\end{thought}

最后，把结论交还现实：本节只建立功能层关系，不由此推出特定脑区实现、主观意识或宇宙目的。可测状态、干预后果和明确失败条件，才是从原文直觉走向可检验理论的桥。

\section{Control Instability}
\begin{narration}
现在把镜头移到“Control Instability”。我们只沿本章已经取得的概念向前一步：先确定它在闭环里的角色，再检查它的尺度与失败边界。
\end{narration}

本节把“Control Instability”落实为把原始身体量转换为强度、导数和置信度组成的主体影响场。这个定义不是说“任何相似现象都算”，而是要求我们同时给出对象域、允许的干预以及观察窗口。对于主题“FCS：现实如何影响主体”而言，它承担的是认知内核通过身体进入现实中的一个可辨识环节。

把这一关系写成最小数学骨架，可得

\begin{equation}f_k=(\sigma_k,\dot\sigma_k,q_k).\end{equation}

式中的符号只承担结构角色，具体参数仍需由任务和身体数据辨识。我们首先检查可观测量，再检查控制量，最后检查比它更慢的约束。这样的顺序来自具身认知、分层控制、MPC、状态估计与安全包络，但本书对Control Instability的命名和组合仍是需要实验验证的建模提案。

这一小节的反例同样重要：原始传感值与功能影响混同会让不同身体无法使用统一语义。因此评价它不能只看一次任务结果，而要比较干预前后的轨迹、恢复时间、维护成本和风险。如果改变该机制而所有允许观察下的分布都不变，那么它至少在当前实验族中是冗余描述。

\begin{thought}
对“Control Instability”做一次消融：冻结它、删除它，或只允许它以相邻时间尺度交换残差。哪些现象立即消失，哪些只在长期出现？若答案无法区分三种干预，我们还没有找到正确的状态变量。
\end{thought}

最后，把结论交还现实：本节只建立功能层关系，不由此推出特定脑区实现、主观意识或宇宙目的。可测状态、干预后果和明确失败条件，才是从原文直觉走向可检验理论的桥。

\section{Urgency}
\begin{narration}
现在把镜头移到“Urgency”。我们只沿本章已经取得的概念向前一步：先确定它在闭环里的角色，再检查它的尺度与失败边界。
\end{narration}

本节把“Urgency”落实为把原始身体量转换为强度、导数和置信度组成的主体影响场。这个定义不是说“任何相似现象都算”，而是要求我们同时给出对象域、允许的干预以及观察窗口。对于主题“FCS：现实如何影响主体”而言，它承担的是认知内核通过身体进入现实中的一个可辨识环节。

把这一关系写成最小数学骨架，可得

\begin{equation}f_k=(\sigma_k,\dot\sigma_k,q_k).\end{equation}

式中的符号只承担结构角色，具体参数仍需由任务和身体数据辨识。我们首先检查可观测量，再检查控制量，最后检查比它更慢的约束。这样的顺序来自具身认知、分层控制、MPC、状态估计与安全包络，但本书对Urgency的命名和组合仍是需要实验验证的建模提案。

这一小节的反例同样重要：原始传感值与功能影响混同会让不同身体无法使用统一语义。因此评价它不能只看一次任务结果，而要比较干预前后的轨迹、恢复时间、维护成本和风险。如果改变该机制而所有允许观察下的分布都不变，那么它至少在当前实验族中是冗余描述。

\begin{thought}
对“Urgency”做一次消融：冻结它、删除它，或只允许它以相邻时间尺度交换残差。哪些现象立即消失，哪些只在长期出现？若答案无法区分三种干预，我们还没有找到正确的状态变量。
\end{thought}

最后，把结论交还现实：本节只建立功能层关系，不由此推出特定脑区实现、主观意识或宇宙目的。可测状态、干预后果和明确失败条件，才是从原文直觉走向可检验理论的桥。

\section{Intensity—Derivative—Confidence}
\begin{narration}
现在把镜头移到“Intensity—Derivative—Confidence”。我们只沿本章已经取得的概念向前一步：先确定它在闭环里的角色，再检查它的尺度与失败边界。
\end{narration}

本节把“Intensity—Derivative—Confidence”落实为把原始身体量转换为强度、导数和置信度组成的主体影响场。这个定义不是说“任何相似现象都算”，而是要求我们同时给出对象域、允许的干预以及观察窗口。对于主题“FCS：现实如何影响主体”而言，它承担的是认知内核通过身体进入现实中的一个可辨识环节。

把这一关系写成最小数学骨架，可得

\begin{equation}f_k=(\sigma_k,\dot\sigma_k,q_k).\end{equation}

式中的符号只承担结构角色，具体参数仍需由任务和身体数据辨识。我们首先检查可观测量，再检查控制量，最后检查比它更慢的约束。这样的顺序来自具身认知、分层控制、MPC、状态估计与安全包络，但本书对Intensity—Derivative—Confidence的命名和组合仍是需要实验验证的建模提案。

这一小节的反例同样重要：原始传感值与功能影响混同会让不同身体无法使用统一语义。因此评价它不能只看一次任务结果，而要比较干预前后的轨迹、恢复时间、维护成本和风险。如果改变该机制而所有允许观察下的分布都不变，那么它至少在当前实验族中是冗余描述。

\begin{thought}
对“Intensity—Derivative—Confidence”做一次消融：冻结它、删除它，或只允许它以相邻时间尺度交换残差。哪些现象立即消失，哪些只在长期出现？若答案无法区分三种干预，我们还没有找到正确的状态变量。
\end{thought}

最后，把结论交还现实：本节只建立功能层关系，不由此推出特定脑区实现、主观意识或宇宙目的。可测状态、干预后果和明确失败条件，才是从原文直觉走向可检验理论的桥。

\section{FCS Geometry $\neq$ SGC Geometry}
\begin{narration}
现在把镜头移到“FCS Geometry $\neq$ SGC Geometry”。我们只沿本章已经取得的概念向前一步：先确定它在闭环里的角色，再检查它的尺度与失败边界。
\end{narration}

本节把“FCS Geometry $\neq$ SGC Geometry”落实为以身体为参考、带持续身份和认识状态的可行动现实关系图。这个定义不是说“任何相似现象都算”，而是要求我们同时给出对象域、允许的干预以及观察窗口。对于主题“FCS：现实如何影响主体”而言，它承担的是认知内核通过身体进入现实中的一个可辨识环节。

把这一关系写成最小数学骨架，可得

\begin{equation}\mathcal G_S=(V,E,X,\eta),\quad \eta\in\{o,d,p,c\}.\end{equation}

式中的符号只承担结构角色，具体参数仍需由任务和身体数据辨识。我们首先检查可观测量，再检查控制量，最后检查比它更慢的约束。这样的顺序来自具身认知、分层控制、MPC、状态估计与安全包络，但本书对FCS Geometry $\neq$ SGC Geometry的命名和组合仍是需要实验验证的建模提案。

这一小节的反例同样重要：预测节点若覆盖观察节点，内部世界会失去现实校准锚点。因此评价它不能只看一次任务结果，而要比较干预前后的轨迹、恢复时间、维护成本和风险。如果改变该机制而所有允许观察下的分布都不变，那么它至少在当前实验族中是冗余描述。

\begin{thought}
对“FCS Geometry $\neq$ SGC Geometry”做一次消融：冻结它、删除它，或只允许它以相邻时间尺度交换残差。哪些现象立即消失，哪些只在长期出现？若答案无法区分三种干预，我们还没有找到正确的状态变量。
\end{thought}

最后，把结论交还现实：本节只建立功能层关系，不由此推出特定脑区实现、主观意识或宇宙目的。可测状态、干预后果和明确失败条件，才是从原文直觉走向可检验理论的桥。

\section*{本章收束：把名词还原为闭环}
现在回头看“FCS：现实如何影响主体”，最重要的不是记住缩写，而是说清本章对象怎样改变系统的状态、代价或可恢复性。下一章“SGC—FCS 二元现实模型”只使用这里已经给出的结果，不反过来为本章提供前提。

\begin{bookproposition}
若一个候选机制既不改变可观测轨迹分布，也不改变资源、风险或恢复代价，那么在当前观察尺度上，它不是一个可辨识的功能结构。
\end{bookproposition}
\begin{proofidea}
功能等价由可干预后果定义。若所有允许干预下的输出分布与代价均不变，则该机制相对于当前实验族不可辨识；要么扩大观察尺度，要么删除这一冗余描述。
\end{proofidea}

\begin{readercheck}
请尝试回答：本章对象的状态是什么？它的最快和最慢时间常数是什么？什么观测能证伪它？它失败时，残差应向哪一层升级？
\end{readercheck}
\cleardoublepage
